\chapter{Spiegabilit\`a dei modelli}
\label{app:xai_complete_review}

In questa appendice sono riportati i risultati completi dell'analisi di spiegabilit\`a condotta tramite SHAP sui dati di test della configurazione arricchita (\texttt{human}).

Per ciascuna strategia di bucketing (\texttt{Prefix Length}, \texttt{Component} e \texttt{Component Time Warped}) e per ciascuno dei quattro classificatori, i risultati sono articolati su tre pagine:
\begin{itemize}
    \item la prima pagina illustra l'attribuzione di importanza nel corso dell'esecuzione a livello Macro, Semantico e Base Feature, affiancata dalla matrice di attenzione temporale complessiva;
    \item la seconda pagina presenta il confronto di fedelt\`a tra il metodo di spiegazione \textit{model-specific} e il benchmark agnostico (\texttt{KernelExplainer});
    \item la terza pagina riporta l'impatto SHAP netto nel corso del tempo sui singoli punteggi delle feature ergonomiche RULA (\texttt{RULA Wrist}, \texttt{RULA Upper Arm}, \texttt{RULA Lower Arm}, \texttt{RULA Neck}).
\end{itemize}
Tutti i grafici sono disposti a tutta larghezza su linea singola per garantire la massima leggibilit\`a e risoluzione lungo l'intera dinamica temporale.

\clearpage

\foreach \strat/\stratTitle/\sShort/\sKey in {
        PREFIX_LENGTH/{Strategia \texttt{Prefix Length}}/{Prefix Length}/pl,
        COMPONENT/{Strategia \texttt{Component}}/{Component}/comp,
        COMPONENT_TIME_WARPED/{Strategia \texttt{Component Time Warped}}/{Time Warped}/ctw
    } {
        \section{\stratTitle}
        \label{sec:app_\sKey}

        \foreach \clf/\clfName in {
                rf/{Random Forest},
                xgb/{XGBoost},
                dt/{Decision Tree},
                lstm/{LSTM}
            } {
                \subsection{Classificatore: \clfName}
                \label{subsec:app_\sKey_\clf}

                % ======================================================================
                % Pagina 1: Dinamica Longitudinale e Attenzione Temporale (4 grafici)
                % ======================================================================
                \begin{figure}[H]
                    \centering
                    % Riga 1: Macro
                    \begin{subfigure}[c]{\textwidth}
                        \centering
                        \includegraphics[width=\textwidth]{images/plots/\strat_\clf_longitudinal_macro.pdf}
                        \vspace{-0.4em}
                        \caption{Distribuzione Macro-Cluster}
                        \label{fig:\sKey_\clf_macro}
                    \end{subfigure}

                    \vspace{0.7em}

                    % Riga 2: Semantico
                    \begin{subfigure}[c]{\textwidth}
                        \centering
                        \includegraphics[width=\textwidth]{images/plots/\strat_\clf_longitudinal_semantic.pdf}
                        \vspace{-0.4em}
                        \caption{Distribuzione Cluster Semantici}
                        \label{fig:\sKey_\clf_sem}
                    \end{subfigure}

                    \vspace{0.7em}

                    % Riga 3: Base Feature
                    \begin{subfigure}[c]{\textwidth}
                        \centering
                        \includegraphics[width=\textwidth]{images/plots/\strat_\clf_longitudinal_base_feature.pdf}
                        \vspace{-0.4em}
                        \caption{Attribuzione Base Feature}
                        \label{fig:\sKey_\clf_base}
                    \end{subfigure}

                    \vspace{0.7em}

                    % Riga 4: Attenzione Temporale
                    \begin{subfigure}[c]{\textwidth}
                        \centering
                        \includegraphics[width=\textwidth]{images/plots/\strat_\clf_temporal_attention.pdf}
                        \vspace{-0.4em}
                        \caption{Matrice Attenzione Temporale ($0\text{--}100\%$)}
                        \label{fig:\sKey_\clf_att}
                    \end{subfigure}
                \end{figure}

                \clearpage

                % ======================================================================
                % Pagina 2: Fedeltà Explainer Model-Specific vs. Model-Agnostic (4 grafici)
                % ======================================================================
                \begin{figure}[H]
                    \centering
                    % Riga 1: Fedeltà Macro
                    \begin{subfigure}[c]{\textwidth}
                        \centering
                        \includegraphics[width=\textwidth]{images/plots/\strat_\clf_top_features_macro.pdf}
                        \vspace{-0.4em}
                        \caption{Fedeltà Macro-Cluster ($100\%$)}
                        \label{fig:\sKey_\clf_fid_macro}
                    \end{subfigure}

                    \vspace{0.4em}

                    % Riga 2: Fedeltà Semantico
                    \begin{subfigure}[c]{\textwidth}
                        \centering
                        \includegraphics[width=\textwidth]{images/plots/\strat_\clf_top_features_semantic.pdf}
                        \vspace{-0.4em}
                        \caption{Fedeltà Cluster Semantici ($100\%$)}
                        \label{fig:\sKey_\clf_fid_sem}
                    \end{subfigure}

                    \vspace{0.4em}

                    % Riga 3: Fedeltà Base Feature
                    \begin{subfigure}[c]{\textwidth}
                        \centering
                        \includegraphics[width=\textwidth]{images/plots/\strat_\clf_top_features_base_feature.pdf}
                        \vspace{-0.4em}
                        \caption{Fedeltà Base Feature ($100\%$)}
                        \label{fig:\sKey_\clf_fid_base}
                    \end{subfigure}

                    \vspace{0.4em}

                    % Riga 4: Fedeltà Atomiche
                    \begin{subfigure}[c]{\textwidth}
                        \centering
                        \includegraphics[width=\textwidth]{images/plots/\strat_\clf_top_features_none.pdf}
                        \vspace{-0.4em}
                        \caption{Fedeltà Feature Atomiche (Top-8)}
                        \label{fig:\sKey_\clf_fid_none}
                    \end{subfigure}
                \end{figure}

                \clearpage

                % ======================================================================
                % Pagina 3: Dinamica Temporale dei Punteggi Ergonomici RULA (4 grafici)
                % ======================================================================
                \begin{figure}[H]
                    \centering
                    % Riga 1: Wrist
                    \begin{subfigure}[c]{\textwidth}
                        \centering
                        \includegraphics[width=\textwidth]{images/plots/\strat_\clf_rula_wrist_heatmap.pdf}
                        \vspace{-0.4em}
                        \caption{Impatto SHAP netto \texttt{RULA Wrist}}
                        \label{fig:\sKey_\clf_rula_wrist}
                    \end{subfigure}

                    \vspace{0.7em}

                    % Riga 2: Upper Arm
                    \begin{subfigure}[c]{\textwidth}
                        \centering
                        \includegraphics[width=\textwidth]{images/plots/\strat_\clf_rula_upper_heatmap.pdf}
                        \vspace{-0.4em}
                        \caption{Impatto SHAP netto \texttt{RULA Upper Arm}}
                        \label{fig:\sKey_\clf_rula_upper}
                    \end{subfigure}

                    \vspace{0.7em}

                    % Riga 3: Lower Arm
                    \begin{subfigure}[c]{\textwidth}
                        \centering
                        \includegraphics[width=\textwidth]{images/plots/\strat_\clf_rula_lower_heatmap.pdf}
                        \vspace{-0.4em}
                        \caption{Impatto SHAP netto \texttt{RULA Lower Arm}}
                        \label{fig:\sKey_\clf_rula_lower}
                    \end{subfigure}

                    \vspace{0.7em}

                    % Riga 4: Neck
                    \begin{subfigure}[c]{\textwidth}
                        \centering
                        \includegraphics[width=\textwidth]{images/plots/\strat_\clf_rula_neck_heatmap.pdf}
                        \vspace{-0.4em}
                        \caption{Impatto SHAP netto \texttt{RULA Neck}}
                        \label{fig:\sKey_\clf_rula_neck}
                    \end{subfigure}
                \end{figure}

                \clearpage
            }
    }
